\chapter{Introdução}
\label{chap:introducao}

A sustentabilidade da dívida pública estadual permanece como uma questão central do federalismo fiscal brasileiro. Desde a década de 1990, diversos estados passaram por episódios de elevado endividamento, renegociações com a União, fragilidade financeira de bancos estaduais e dificuldades recorrentes de ajuste fiscal. Esse histórico revelou uma tensão importante do arranjo federativo nacional: embora os governos estaduais tenham autonomia para executar políticas públicas e assumir compromissos financeiros, sua dinâmica fiscal é condicionada por regras federais, decisões macroeconômicas nacionais e pela expectativa de renegociação em momentos de crise. Nesse contexto, problemas de risco moral, restrição orçamentária branda e uso comum de recursos fiscais tornam-se relevantes para entender os incentivos ao endividamento subnacional, como discutido por \citeonline{velasco2000} e \citeonline{rodden2002}.

A crise fiscal dos anos 1990 levou à construção de mecanismos institucionais voltados a disciplinar o comportamento fiscal dos estados. Entre esses mecanismos, destaca-se a Lei n\textsuperscript{o}~9.496/1997, que autorizou a União a assumir e refinanciar dívidas estaduais mediante contratos individuais com os entes federativos. Esses contratos não apenas reorganizaram o perfil financeiro dos passivos, mas também impuseram condições de ajuste, prazos de pagamento, encargos e limites para o comprometimento da Receita Corrente Líquida (RCL) com o serviço da dívida. Assim, mais do que uma simples operação de refinanciamento, a Lei n\textsuperscript{o}~9.496/1997 estruturou um arranjo contratual de disciplina fiscal subnacional, no qual o alívio financeiro concedido pela União foi acompanhado por restrições ao comportamento fiscal futuro dos estados.

A Lei de Responsabilidade Fiscal (LRF), aprovada em 2000, ampliou essa lógica em um marco institucional mais abrangente. Enquanto a Lei n\textsuperscript{o}~9.496/1997 operava por meio de contratos bilaterais entre União e estados refinanciados, a LRF estabeleceu normas gerais de finanças públicas para os diferentes níveis de governo, com ênfase em planejamento, transparência, limites de despesa, controle do endividamento e responsabilização dos gestores. Portanto, a LRF não deve ser interpretada como simples continuação automática da Lei n\textsuperscript{o}~9.496/1997. A relação entre ambas é mais bem compreendida como uma sequência institucional: a renegociação de 1997 introduziu disciplina fiscal em bases contratuais no contexto da crise estadual, enquanto a LRF consolidou um regime mais amplo de responsabilidade fiscal para o setor público brasileiro.

Ainda assim, a recorrência de novas renegociações e regimes especiais de recuperação fiscal sugere que o problema da dívida estadual não foi plenamente resolvido. Esse ponto não implica afirmar, de antemão, que as regras fiscais foram ineficazes, mas indica que sua efetividade deve ser investigada empiricamente. Em particular, é necessário avaliar se os estados brasileiros reagem ao aumento do endividamento por meio de maior esforço fiscal e se essa reação varia conforme o desenho institucional dos contratos de refinanciamento. A questão central, portanto, não é apenas saber se existem regras fiscais, mas se o modo como essas regras são desenhadas e aplicadas altera os incentivos dos governos estaduais.

A literatura sobre regras fiscais destaca justamente que sua efetividade depende menos da existência formal da regra e mais de sua qualidade institucional, de seus mecanismos de \emph{enforcement} e de sua capacidade de lidar com incentivos políticos adversos. Regras fiscais podem reduzir o viés deficitário e fortalecer a credibilidade da política fiscal, como mostram \citeonline{poterba1994} e \citeonline{badinger2017}, mas também podem produzir efeitos limitados quando são mal calibradas, pouco críveis ou aplicadas em contextos de fraca disciplina institucional, como discutido por \citeonline{heinemann2018}, \citeonline{caselli2020} e \citeonline{brandleFiscalRulesMatter2024}. No caso subnacional, esse debate é particularmente relevante porque as decisões fiscais dos estados ocorrem em ambiente federativo, no qual a possibilidade de apoio do governo central pode alterar os incentivos ao ajuste.

Do ponto de vista da sustentabilidade da dívida, este trabalho parte da abordagem de \citeonline{bohn1998} e \citeonline{bohn2007}. A literatura inicial avaliava a solvência fiscal principalmente por meio de testes de estacionariedade e cointegração, associados à restrição orçamentária intertemporal do governo, como em \citeonline{hamilton1985} e \citeonline{trehan1988}. Bohn propõe uma alternativa baseada na função de reação fiscal: se o resultado primário aumenta sistematicamente quando a dívida cresce, há evidência de que a política fiscal contém mecanismos de estabilização do endividamento. Essa abordagem é útil porque desloca o foco da trajetória isolada da dívida para o comportamento fiscal do governo diante do aumento dos passivos. No entanto, essa interpretação exige cautela: uma resposta positiva do resultado primário à dívida deve ser entendida como evidência compatível com sustentabilidade fiscal, e não como prova definitiva de solvência em todos os estados ou em todos os cenários.

Este TCC aplica e adapta essa lógica ao caso dos estados brasileiros. O ponto de partida metodológico é o estudo de \citeonline{abubakarDebtReliefFiscal2025}, que investiga se regras fiscais e programas de alívio da dívida melhoram a sustentabilidade fiscal em países da África Subsaariana. A estratégia original combina dois modelos complementares: uma função de reação fiscal, voltada a avaliar se o resultado primário responde ao estoque de dívida, e um modelo de dinâmica da dívida, destinado a investigar os fatores associados à evolução do endividamento. A presente pesquisa adapta essa estrutura para o contexto subnacional brasileiro, preservando a lógica empírica central, mas substituindo variáveis sem equivalente direto no caso dos estados.

A primeira adaptação diz respeito ao denominador fiscal. Enquanto grande parte da literatura internacional trabalha com dívida e resultado primário em proporção do PIB, este trabalho utiliza a Receita Corrente Líquida como referência principal. Essa escolha decorre da relevância institucional da RCL na LRF e nos contratos de refinanciamento estadual, além de sua maior aderência às regras fiscais aplicáveis aos estados brasileiros. Assim, a sustentabilidade fiscal é analisada a partir da Dívida Consolidada Líquida sobre RCL (DCL/RCL) e do resultado primário sobre RCL, aproximando a mensuração empírica do arcabouço legal efetivamente enfrentado pelos governos estaduais.

A segunda adaptação refere-se à mensuração das regras fiscais. Em vez de utilizar apenas uma variável binária de existência de regra fiscal, o trabalho explora a heterogeneidade dos limites contratuais de comprometimento da RCL definidos nos contratos da Lei n\textsuperscript{o}~9.496/1997. Esses limites variam entre estados e funcionam como regras fiscais contratuais, pois restringem o serviço da dívida em relação à capacidade de pagamento medida pela receita. Em alguns contratos, o limite foi definido como intervalo, e não como percentual único; nesses casos, utiliza-se o limite superior como medida principal. Dessa forma, a variável de teto deve ser interpretada como o limite máximo contratual de comprometimento da RCL, isto é, uma medida de permissividade contratual máxima.

A pergunta de pesquisa pode ser formulada em duas dimensões. Primeiro, os estados brasileiros apresentam uma reação fiscal sistemática ao aumento do endividamento, compatível com a abordagem de Bohn? Segundo, essa reação é condicionada pelo limite contratual de comprometimento da RCL estabelecido nos contratos de refinanciamento da Lei n\textsuperscript{o}~9.496/1997? Complementarmente, o trabalho investiga a dinâmica da DCL/RCL, buscando avaliar em que medida a persistência da dívida, o resultado primário e as condições econômicas estão associados à trajetória do endividamento estadual.

Para responder a essas questões, utiliza-se um painel de estados brasileiros no período de 2002 a 2024. O Modelo I estima uma função de reação fiscal em que o resultado primário sobre RCL depende da DCL/RCL defasada, da interação entre dívida defasada e teto contratual, do ciclo econômico estadual e de efeitos fixos de estado e ano. A estimação é realizada por mínimos quadrados em dois estágios (2SLS), uma vez que tanto a dívida defasada quanto sua interação com o teto podem estar sujeitas a endogeneidade. O Modelo II, por sua vez, examina a dinâmica da dívida estadual, modelando a DCL/RCL como função de sua própria defasagem, do resultado primário, do crescimento econômico e de controles adicionais. Como se trata de um modelo dinâmico em painel, são empregados estimadores apropriados para lidar com persistência e possível viés de pequena amostra, incluindo LSDVC e especificações de robustez.

A contribuição do trabalho é dupla. Em primeiro lugar, aplica-se ao caso brasileiro uma estratégia empírica inspirada em \citeonline{abubakarDebtReliefFiscal2025}, adaptando-a às instituições fiscais estaduais e à disponibilidade de dados subnacionais. Em segundo lugar, o trabalho desloca o foco da simples existência de regras fiscais para o desenho institucional dessas regras, explorando a variação dos tetos contratuais da Lei n\textsuperscript{o}~9.496/1997 como elemento capaz de condicionar a reação fiscal dos estados ao endividamento. Dessa forma, a análise busca contribuir para o debate sobre sustentabilidade fiscal subnacional ao investigar não apenas se há ajuste fiscal diante do aumento da dívida, mas também se esse ajuste varia conforme a permissividade dos contratos de refinanciamento.

O restante do trabalho está organizado da seguinte forma. A Seção~\ref{chap:revisao} revisa a literatura sobre regras fiscais, federalismo fiscal e sustentabilidade da dívida. A Seção~\ref{chap:institucional} apresenta o contexto institucional brasileiro, com ênfase na Lei n\textsuperscript{o}~9.496/1997, na LRF e nas renegociações da dívida estadual. A Seção~\ref{chap:metodologia} descreve os dados, as variáveis e a estratégia empírica. A Seção~\ref{chap:resultados} apresenta e discute os resultados dos modelos de reação fiscal e dinâmica da dívida. Por fim, a Seção~\ref{chap:conclusao} reúne as principais conclusões, limitações e implicações para o desenho de regras fiscais subnacionais.